% ===================================================================
%  TABELO N-RO 4 — Matura aĝo kaj maljuneco.
%  Traduction 2026 d'apres texto/io, controlee sur texto/fr et
%  texto/es. Consigne : voir texto/eo/10-tabelo-01.tex.
% ===================================================================

%%K t04-apar-1 apar
\VUsaut{4.0mm}
\VUtitre{15.0pt}{60}{TABELO N\textsuperscript{o} 4}
\VUsaut{1.6mm}
\VUtitre{11.4pt}{0}{Matura aĝo kaj maljuneco.}
\VUsaut{3.2mm}

%%K t04-tit-1 sub
\VUcentre{11.4pt}{0}{\textit{Unua sceno.}}
\VUsaut{1.6mm}
\VUtitre{11.4pt}{0}{La edziĝa festeno.}
\VUsaut{2.6mm}

%%K t04-tit-2 sub
\VUpk{10.2pt}{40}{La aŭtuno.}
\VUsaut{3.0mm}

%%K t04-c1-01-1 p
1. --- La junuloj plikreskis. Unu el ili, Johano, edziĝas. Ni ĉeestas la
\VUgras{edziĝan festenon} kiu kunigas, ĉirkaŭ la sama tablo, la familiojn de
la geedzoj.

%%K t04-c1-02-1 p
2. --- Ni estas en \VUgras{aŭtuno}, en la \VUgras{monato septembro}. La
malvarma vento de \VUgras{oktobro} kaj de \VUgras{novembro} ankoraŭ ne
senigis la kamparon je ĝia verda foliaro. La vespera aero estas tepida kaj
parfuma; pro tio la vespermanĝo okazas sub la \VUgras{verando}
\textsuperscript{(8)} ĉe la ĝardeno.

%%K t04-c1-03-1 p
3. --- Malproksime, sur la \VUgras{monteto} \textsuperscript{(3)}, situas la
hejmo de la gepatroj de Johano, eleganta \VUgras{kastelo} \textsuperscript{(1)}
kaŝita en la verdaĵo; belaj vitejoj, kiuj produktas bonegan vinon, kovras la
deklivon de la monteto. La vitkultivisto kultivas kaj zorgas ilin.

%%K t04-c1-04-1 p
4. --- Super la vitejoj, pasas flugo de \VUgras{perdrikoj}
\textsuperscript{(2)} timigitaj de la alproksimiĝo de la \VUgras{ĉasgardisto}
\textsuperscript{(5)}. Ĉi tiu, kun \VUgras{pafilo} sub la brako, kaj
\VUgras{ŝultra ĉassako}, esploras la \VUgras{arbustaron}
\textsuperscript{(4)} kun sia \VUgras{hundo} \textsuperscript{(6)}. Li
gvatas la bienon kaj malhelpas la ŝtelĉasistojn detrui la ĉasaĵon.

%%K t04-c1-05-1 p
5. --- Ekster la \VUgras{verando}, grimpas \VUgras{vitespaliro} de kiu pendas
densaj \VUgras{vinberaroj} \textsuperscript{(7)}.

%%K t04-c1-06-1 p
6. --- La belaj aŭtunaj floroj, kiuj garnas la \VUgras{tablon}
\textsuperscript{(16)}, estas \VUgras{krizantemoj} \textsuperscript{(9)}; la
\VUgras{patro} \textsuperscript{(10)} de la edziniĝinto metis unu el ili en la
butontruon de sia frako.

%%K t04-c1-07-1 p
7. --- La manĝo finiĝas. La \VUgras{servisto} \textsuperscript{(11)} komencas
alporti la desertpladojn kaj baldaŭ deprenos la diversajn partojn de la
manĝilaro, \VUgras{kulerojn} \textsuperscript{(14)}, \VUgras{forketojn}
\textsuperscript{(12)}, \VUgras{tranĉilojn} \textsuperscript{(13)},
\VUgras{pladegojn} \textsuperscript{(15)} kiujn oni ne plu bezonas.

%%K t04-tit-3 sub
\VUpk{10.2pt}{40}{La parencaro.}
\VUsaut{3.0mm}

%%K t04-c1-08-1 p
8. --- Ĉiuj aspektas ĝojaj, kaj la rideto per kiu la \VUgras{patrino}
\textsuperscript{(17)} de la edziĝinto rigardas siajn filojn pruvas ŝian
feliĉon.

%%K t04-c1-09-1 p
9. --- Tiu edziĝo unuigas du riĉajn familiojn de la lando. Johano, la
\VUgras{edzo} \textsuperscript{(18)}, estas \VUgras{filo} de granda
bienposedanto kaj li fariĝas la \VUgras{bofilo} de lerta industriisto.

%%K t04-c1-10-1 p
10. --- La \VUgras{geedzoj} estas junaj kaj tamen ili delonge estis
\VUgras{fianĉigitaj}. La \VUgras{juna virino} \textsuperscript{(19)}, Marta,
estas ĉarma en sia \VUgras{blanka robo} garnita per floroj de oranĝarbo.

%%K t04-c1-11-1 p
11. --- Ŝia \VUgras{bopatro} \textsuperscript{(20)} estas tre feliĉa havi tiel
afablan \VUgras{bofilinon}, kaj la \VUgras{bopatrino} \textsuperscript{(21)} de
Johano ridetas vidante tiel bela sian \VUgras{filinon}.

%%K t04-c1-12-1 p
12. --- Ĉe la ekstremo de la tablo, sidas la ceteraj \VUgras{invititoj}
\textsuperscript{(22)}, \VUgras{parencoj, amikoj, kuzoj, kuzinoj}.
\VUgras{Tablestro} \textsuperscript{(23)}, kun servotuko sub la brako,
diligente servas ilin.

%%K t04-tit-4 sub
\VUpk{10.2pt}{40}{La amikoj.}
\VUsaut{3.0mm}

%%K t04-c1-13-1 p
13. --- Dekstraflanke de la tablo, jen \VUgras{fraŭlino} \textsuperscript{(24)},
\VUgras{amikino} de la juna edzino, poste la \VUgras{frato} de ĉi tiu, kiu
estas ŝia \VUgras{honorkunulo} \textsuperscript{(25)}. Li babilas kun sia
\VUgras{honorfraŭlino} \textsuperscript{(26)} fratino de la edzo kaj kiu, per
tiu edziĝo, fariĝas la \VUgras{bofratino} de Marta. Ŝi estas ĉarma junulino.
Ŝi havas belajn \VUgras{juvelojn, perlokolĉenon}, kaj oran \VUgras{braceleton}.

%%K t04-c1-14-1 p
14. --- Apud ili, la sinjoro kiu staras kaj levas sian glason estas
\VUgras{amiko} \textsuperscript{(27)} de la familio. Li estas unu el la
\VUgras{atestantoj de la edzo}. Li \VUgras{tostas}, trinkas por la
\VUgras{sano} de la junaj geedzoj kaj deziras al ili longan feliĉon. Li estas
vestita per ceremonia kostumo, per frako kun \VUgras{lakitaj ŝuoj}
\textsuperscript{(28)}. Li estas la akompananto de la \VUgras{avino}
\textsuperscript{(29)} kiu estas \VUgras{vidvino}.

%%K t04-c1-15-1 p
15. --- La junulo en \VUgras{frako} \textsuperscript{(31)}, kun maldika nigra
lipbarbo, estas la \VUgras{bofrato} \textsuperscript{(30)} de Johano. Li estas
eminenta inĝeniero. Li estas la najbaro de \VUgras{juna sinjorino}
\textsuperscript{(33)} tre eleganta, la \VUgras{pli aĝa fratino} de Marta kaj
la patrino de la ĉarma knabineto kiu venas, donante la brakon al sia kuzo,
oferi floron kaj \VUgras{deziri} al li \VUgras{bonan vesperon}. Tiu sinjorino
havas tre belajn \VUgras{orelringojn} kaj elegantan \VUgras{kapveston}.

%%K t04-c1-16-1 p
16. --- La malgranda kuzo, tiel stranga pro sia \VUgras{bluzo}
\textsuperscript{(36)} kaj sia pufa \VUgras{pantalono} \textsuperscript{(37)}
estas tre kompatinda. Li estas \VUgras{orfo} \textsuperscript{(35)}. Tre juna
li perdis sian patron kaj sian patrinon. Lia onklo estas lia \VUgras{tutoro}
\textsuperscript{(38)}, kaj restis fraŭlo por eduki la kompatindan infanon. Li
estas bonkora viro. Li estas iom kalva kaj tre miopa; li uzas
\VUgras{(naz)binoklon}.

%%K t04-tit-5 sub
\VUsaut{2.4mm}
\VUpk{10.2pt}{40}{La deserto.}
\VUsaut{3.0mm}

%%K t04-c1-17-1 p
17. --- Malantaŭ la festenanoj, sur tablo kovrita per \VUgras{(tablo)tuko}, la
\VUgras{deserto} \textsuperscript{(39)} estas preparita. En pelvego, jen
fruktoj, \VUgras{persikoj} \textsuperscript{(40)}, \VUgras{piroj}
\textsuperscript{(41)}, \VUgras{pomoj} \textsuperscript{(42)},
\VUgras{abrikotoj} \textsuperscript{(43)} kaj \VUgras{vinberoj}
\textsuperscript{(44)}. Proksime, \VUgras{ananasoj} \textsuperscript{(45)},
bela \VUgras{kuko} \textsuperscript{(46)}, \VUgras{likvorboteloj}
\textsuperscript{(48)} kaj \VUgras{ĉampano} \textsuperscript{(49)} kaj ankaŭ
kolonoj da \VUgras{teleroj} \textsuperscript{(47)} puraj.

%%K t04-c1-18-1 p
18. --- La \VUgras{kelisto} \textsuperscript{(50)} malŝtopas \VUgras{botelon}
\textsuperscript{(52)} de malnova vino per \VUgras{korktirilo}
\textsuperscript{(51)}. La \VUgras{korka ŝtopilo} \textsuperscript{(53)}
ŝajnas tre malfacile eltirebla. Tiu kelisto havas \VUgras{servotukon}
\textsuperscript{(54)} sub la brako.

%%K t04-c1-19-1 p
19. --- Post la vespermanĝo, la sinjoroj iros al la fumejo, kaj la sinjorinoj
iros pretigi sin por la balo.

%%K t04-tit-6 sub
\VUsaut{5.0mm}
\VUcentre{11.4pt}{0}{\textit{Dua sceno.}}
\VUsaut{1.6mm}
\VUpk{11.4pt}{40}{La datreveno de la avo.}
\VUsaut{2.6mm}

%%K t04-tit-7 sub
\VUpk{10.2pt}{40}{La vizito.}
\VUsaut{3.0mm}

%%K t04-c2-01-1 p
1. --- Dek jaroj pasis, la gepatroj de Johano fariĝis maljunaj tre amantaj
siajn tri genepojn, du \VUgras{nepojn} \textsuperscript{(2)} kaj unu
\VUgras{nepinon} \textsuperscript{(3)}. Ĉar hodiaŭ okazas \VUgras{la datreveno
de la avo}, la infanoj venas deziri al li bonan naskiĝfeston.

%%K t04-c2-02-1 p
2. --- La malgranda Petro estas ankoraŭ tre juna, li aĝas nur tri jarojn. Li
vidas la \VUgras{katon} \textsuperscript{(1)} alproksimiĝi. Li volas karesi
ĝian belan \VUgras{felon} silkecan kaj ĝian longan \VUgras{voston}, li ne
pripensas, ke sub la graciaj \VUgras{piedoj}, kaŝiĝas maliceaj ungoj pintaj
kiuj eble gratos lin.

%%K t04-c2-03-1 p
3. --- Lia fratino Gabriela, kiu donas al li la manon, estas multe pli serioza
kaj Marcelo kun sia granda \VUgras{mantelo} \textsuperscript{(4)} kaj leda
\VUgras{zono} \textsuperscript{(5)} aspektas iom vireca, malgraŭ sia mallonga
pantalono kiu lasas vidi siajn \VUgras{ŝtrumpojn} \textsuperscript{(6)}.

%%K t04-c2-04-1 p
4. --- Ilia patro metis sian dikan vintran \VUgras{surtuton}
\textsuperscript{(7)} kun \VUgras{felkolumo} \textsuperscript{(8)} kaj, en sia
\VUgras{poŝo} \textsuperscript{(9)}, li havas koltukon. Lia edzino havas sian
feltan ĉapelon garnitan per \VUgras{plumoj} \textsuperscript{(10)}, sian
\VUgras{jaketon} \textsuperscript{(11)}, sian \VUgras{boa-felaĵon}
\textsuperscript{(12)} kaj sian \VUgras{mufon} \textsuperscript{(13)}.

%%K t04-c2-05-1 p
5. --- Estas tre malvarme, ĉar regas la vintro. La \VUgras{neĝo}
\textsuperscript{(19)} falis dum la tuta tago kaj la tegmentoj de la domoj
estas tute blankaj. Samtempe kiel la \VUgras{nokto}, la malvarmo fariĝas eĉ
pli akra. En la ĉielo, la \VUgras{luno} \textsuperscript{(17)} kaj la
\VUgras{steloj} \textsuperscript{(18)} brilas kaj trapasas \VUgras{la
mallumon}.

%%K t04-tit-8 sub
\VUsaut{4.2mm}
\VUpk{10.2pt}{40}{La malsano.}
\VUsaut{3.0mm}

%%K t04-c2-06-1 p
6. --- La kompatinda \VUgras{blindulo} \textsuperscript{(20)} kiu trairas la
placon antaŭ \VUgras{la apoteko} \textsuperscript{(16)} kondukata de sia
\VUgras{pudelo} \textsuperscript{(21)}, estas tre malfeliĉa. Justina, la
\VUgras{servistino} \textsuperscript{(14)}, alportas el la kuirejo
\VUgras{pelvon} \textsuperscript{(15)} da tre varma \VUgras{tizano}
sukrigita malavare.

%%K t04-c2-07-1 p
7. --- La \VUgras{avino} \textsuperscript{(23)} kiu volas serĉi sian
\VUgras{orelbinoklon} \textsuperscript{(26)} sur la tablo, por legi la
\VUgras{recepton} \textsuperscript{(27)} kiun la \VUgras{kuracisto}
\textsuperscript{(25)} ĵus skribis, sin levas de sia brakseĝo, apogante sin sur
sia \VUgras{bastono} \textsuperscript{(24)}.

%%K t04-c2-08-1 p
8. --- Ofte ŝi suferas pro \VUgras{malsanetoj, kapturnoj, kataretoj,
dentodoloroj}, ŝiaj kruroj estas malfortaj kaj ŝiaj okuloj ne plu estas tre
bonaj. Malgraŭ tio, volonte ŝi mokas sian maljunan \VUgras{kuraciston} kiu
ĉiam volas preskribi al ŝi \VUgras{pocion} \textsuperscript{(28)}. Hodiaŭ ŝi
estas tre ĝoja ĉar ŝi havas ĉirkaŭ si sian junan familion.

%%K t04-c2-09-1 p
9. --- Estas komforte en la ĉambro bone fermita. En la marmora \VUgras{kameno}
\textsuperscript{(29)} flamas \VUgras{belega fajro} \textsuperscript{(30)}, kaj
oni sentas sin feliĉa en la dolĉa \VUgras{lumo} de la \VUgras{lampo}
\textsuperscript{(31)} kiun oni ĵus ekbruligis.

%%K t04-tit-9 sub
\VUsaut{4.2mm}
\VUpk{10.2pt}{40}{La avo.}
\VUsaut{3.0mm}

%%K t04-c2-10-1 p
10. --- Eĉ la \VUgras{avo} \textsuperscript{(33)} forgesis, ke li estis
malsana, ke lia \VUgras{reŭmatismo} fiksas lin ĉe sia \VUgras{brakseĝo}
\textsuperscript{(34)} kaj ke, hieraŭ ankoraŭ, li havis \VUgras{febron kaj
sinkopojn}. Li ne plu memoras, ke dum la lasta vintro li suferis pro
\VUgras{pulmonito} kaj ke raŭka kaj peniga \VUgras{tusado} igis lin tre
suferanta.

%%K t04-c2-10-2 p
Kaj tamen tre malfacile li estis resanigita. De nun li fartas iom pli bone. Sed
lia kruro kiun li apogas sur \VUgras{kuseno} \textsuperscript{(38)} metita sur
malgranda \VUgras{tabureto} \textsuperscript{(39)} ankaŭ doloras lin.

%%K t04-c2-11-1 p
11. --- Kiam li estas zorge envolvita en sia varma \VUgras{ĉambrorobo}
\textsuperscript{(35)} kun dikaj \VUgras{butonoj} \textsuperscript{(36)}
perlamotaj kaj kiam li havas apud si, por legi al li la ĵurnalon, la
malgrandan \VUgras{orfinon} \textsuperscript{(40)} vestitan per blanka
\VUgras{kufo}, per blua \VUgras{robo} \textsuperscript{(42)} kaj
\VUgras{pelerino} \textsuperscript{(41)}, li ne plendas.

%%K t04-c2-12-1 p
12. --- Ĉiujare, kiam la \VUgras{kalendaro} \textsuperscript{(43)} denove
montras la \VUgras{daton} de la unua \VUgras{decembro}, li senpacience atendas
siajn \VUgras{genepojn}, ĉar li scias, ke ili ne forgesos tiun \VUgras{daton}
gravan.

%%K t04-c2-13-1 p
13. --- Li emocie memoras la tempon tre antaŭan kiam li persone venis deziri
bonan feston al la glora imperia \VUgras{marŝalo}, kies \VUgras{portreto}
\textsuperscript{(44)} pendas ĉe la muro, kaj kiu estas la \VUgras{praavo} de
liaj infanoj.
\VUsaut{6.0mm}
\VUfilet{22mm}
